% META: {"id": "TEXP-MATRICES-MARKOV-QCM", "chapitre": "TEXP-MATRICES-MARKOV", "type_objet": "qcm", "genere_depuis": "chapitres/TEXP-MATRICES-MARKOV/qcm/TEXP-MATRICES-MARKOV-QCM.json", "status": "generated"}
% Fichier genere par scripts/build_qcm_tex.py — ne pas editer a la main.

\section*{\textcolor{chapcolor}{\MakeUppercase{Matrices et chaines de Markov}}}

\begin{center}
\textit{Pour chaque question, une seule réponse est exacte.}
\end{center}

\begin{enumerate}

\item[] \textbf{Capacité C4}

\item \textbf{[Q1]} Une matrice de transition d'une chaine de Markov a pour propriété :
  \begin{enumerate}[label=\Alph*.]
    \item La somme de chaque ligne vaut 1
    \item La somme des colonnes vaut 1
    \item Tous ses coefficients sont négatifs
    \item Elle est toujours symetrique
  \end{enumerate}

\item[] \textbf{Capacité C5}

\item \textbf{[Q2]} Si $P$ est la matrice de transition et $P_0$ l'etat initial (vecteur ligne), alors $P_n =$
  \begin{enumerate}[label=\Alph*.]
    \item $P_0 P^n$
    \item $P^n P_0$
    \item $P_0 + n P$
    \item $P_0 P n$
  \end{enumerate}

\item[] \textbf{Capacité C7}

\item \textbf{[Q3]} Un etat stable $X$ d'une chaine de Markov vérifie :
  \begin{enumerate}[label=\Alph*.]
    \item $P X = X$
    \item $X P = X$ et la somme des composantes vaut 1
    \item $X P = 0$
    \item $X^2 = P$
  \end{enumerate}

\item[] \textbf{Capacité C3}

\item \textbf{[Q4]} Une suite de matrices colonnes vérifie $U_{n+1}=AU_n+C$. Une solution constante $U$ satisfait :
  \begin{enumerate}[label=\Alph*.]
    \item $AU=C$
    \item $U=A+C$
    \item $(I-A)U=C$
    \item $U=0$ dans tous les cas
  \end{enumerate}

\item[] \textbf{Capacité C1}

\item \textbf{[Q5]} Pour inverser une matrice 2x2 $A = \begin{pmatrix} a & b \\ c & d \end{pmatrix}$, le determinant est :
  \begin{enumerate}[label=\Alph*.]
    \item $a+b+c+d$
    \item $ac - bd$
    \item $ad + bc$
    \item $ad - bc$
  \end{enumerate}

\item[] \textbf{Capacité C2}

\item \textbf{[Q6]} Une usine compte trois ateliers ; on note $m_{ij}$ le nombre de pièces transférées de l'atelier $i$ vers l'atelier $j$. Modéliser cette situation par une matrice, c'est :
  \begin{enumerate}[label=\Alph*.]
    \item écrire la matrice colonne des totaux produits par chaque atelier
    \item écrire la matrice carrée $3\times3$ dont le coefficient de la ligne $i$ et de la colonne $j$ est $m_{ij}$
    \item écrire une matrice $3\times3$ nécessairement symétrique
    \item écrire la matrice ligne des trois ateliers
  \end{enumerate}

\item[] \textbf{Capacité C6}

\item \textbf{[Q7]} Pour une chaîne de Markov de matrice de transition $T$ et de distribution initiale $\pi_0$ écrite en ligne, la distribution après $n$ transitions vaut :
  \begin{enumerate}[label=\Alph*.]
    \item $T^n\pi_0$
    \item $\pi_0+nT$
    \item $\pi_0T^n$
    \item $n\pi_0T$
  \end{enumerate}

\end{enumerate}

% NEXUS-QCM-TEACHER-ONLY-BEGIN
\ifnxVersionProfesseur
\clearpage
\section*{Clé de correction — réservée au professeur}

\begin{center}
\begin{tabular}{lll}
\hline
Question & Capacité & Réponse exacte \\
\hline
Q1 & C4 & \textbf{A} \\
Q2 & C5 & \textbf{A} \\
Q3 & C7 & \textbf{B} \\
Q4 & C3 & \textbf{C} \\
Q5 & C1 & \textbf{D} \\
Q6 & C2 & \textbf{B} \\
Q7 & C6 & \textbf{C} \\
\hline
\end{tabular}
\end{center}

\subsection*{Diagnostic des réponses erronées}

\begin{itemize}
  \item \textbf{Q1}
  \begin{itemize}
    \item \textbf{B} — Confusions entre lignes et colonnes. Avec la convention des distributions lignes, chaque ligne d'une matrice de transition contient des probabilités dont la somme vaut 1. \emph{Renvoi : C4.}
    \item \textbf{C} — Les probabilités sont positives. Avec la convention des distributions lignes, chaque ligne d'une matrice de transition contient des probabilités dont la somme vaut 1. \emph{Renvoi : C4.}
    \item \textbf{D} — Symetrie non requise. Avec la convention des distributions lignes, chaque ligne d'une matrice de transition contient des probabilités dont la somme vaut 1. \emph{Renvoi : C4.}
  \end{itemize}
  \item \textbf{Q2}
  \begin{itemize}
    \item \textbf{B} — Ordre des facteurs incorrect avec vecteur ligne. Après n transitions, une distribution ligne initiale pi\_0 devient pi\_n=pi\_0 P\textasciicircum{}n. \emph{Renvoi : C5.}
    \item \textbf{C} — Addition au lieu de produit matriciel. Après n transitions, une distribution ligne initiale pi\_0 devient pi\_n=pi\_0 P\textasciicircum{}n. \emph{Renvoi : C5.}
    \item \textbf{D} — Multiplication par n au lieu de la puissance n. Après n transitions, une distribution ligne initiale pi\_0 devient pi\_n=pi\_0 P\textasciicircum{}n. \emph{Renvoi : C5.}
  \end{itemize}
  \item \textbf{Q3}
  \begin{itemize}
    \item \textbf{A} — Vecteur colonne au lieu de vecteur ligne. Une distribution invariante ligne X vérifie XP=X, avec composantes positives de somme 1. \emph{Renvoi : C7.}
    \item \textbf{C} — Produit nul incorrect. Une distribution invariante ligne X vérifie XP=X, avec composantes positives de somme 1. \emph{Renvoi : C7.}
    \item \textbf{D} — Carre matriciel incorrect. Une distribution invariante ligne X vérifie XP=X, avec composantes positives de somme 1. \emph{Renvoi : C7.}
  \end{itemize}
  \item \textbf{Q4}
  \begin{itemize}
    \item \textbf{A} — Substituer une suite constante donne U=AU+C ; oublier le terme U du membre gauche mène à l'équation incomplète AU=C. Une solution constante U de U\_(n+1)=AU\_n+C vérifie U=AU+C, donc (I-A)U=C. \emph{Renvoi : C3.}
    \item \textbf{B} — L'addition A+C mélange une matrice et un vecteur sans appliquer A à U et ne traduit pas la récurrence. Une solution constante U de U\_(n+1)=AU\_n+C vérifie U=AU+C, donc (I-A)U=C. \emph{Renvoi : C3.}
    \item \textbf{D} — Le vecteur nul n'est solution que si C=0 ; aucune telle hypothèse n'est donnée. Une solution constante U de U\_(n+1)=AU\_n+C vérifie U=AU+C, donc (I-A)U=C. \emph{Renvoi : C3.}
  \end{itemize}
  \item \textbf{Q5}
  \begin{itemize}
    \item \textbf{A} — Somme des termes. Le déterminant de la matrice [[a,b],[c,d]] est ad-bc. \emph{Renvoi : C1.}
    \item \textbf{B} — Erreur de formule du determinant. Le déterminant de la matrice [[a,b],[c,d]] est ad-bc. \emph{Renvoi : C1.}
    \item \textbf{C} — Signe + au lieu de -. Le déterminant de la matrice [[a,b],[c,d]] est ad-bc. \emph{Renvoi : C1.}
  \end{itemize}
  \item \textbf{Q6}
  \begin{itemize}
    \item \textbf{A} — Trois totaux ne retiennent pas qui transfère vers qui : l'information de flux, qui est justement l'objet du modèle, est perdue. \emph{Renvoi : C2.}
    \item \textbf{C} — Rien n'impose $m_{ij}=m_{ji}$ : les transferts n'ont aucune raison d'être réciproques. Imposer la symétrie fausserait le modèle. \emph{Renvoi : C2.}
    \item \textbf{D} — Une matrice ligne ne code que trois nombres, alors que la situation en comporte neuf, un par couple d'ateliers. \emph{Renvoi : C2.}
  \end{itemize}
  \item \textbf{Q7}
  \begin{itemize}
    \item \textbf{A} — Avec une distribution écrite en ligne, ce produit n'est même pas défini. La récurrence $\pi_{n+1}=\pi_nT$ se déroule en $\pi_n=\pi_0T^n$ ; c'est la convention colonne qui donnerait $T^n\pi_0$. \emph{Renvoi : C6.}
    \item \textbf{B} — Une distribution et une matrice de transition ne s'additionnent pas : elles ne décrivent pas les mêmes objets, et la somme des coefficients cesserait de valoir $1$. \emph{Renvoi : C6.}
    \item \textbf{D} — Multiplier par $n$ ferait sortir la somme des coefficients de $1$. Le nombre de transitions intervient en \textbf{puissance} de $T$, par itération de $\pi_{n+1}=\pi_nT$. \emph{Renvoi : C6.}
  \end{itemize}
\end{itemize}
\fi
% NEXUS-QCM-TEACHER-ONLY-END
